% !TEX root = main.tex
\section{Judge Prompts and Scoring Configuration}\label{app:judge-prompts}
\paragraph{Judges and prompts of the transfer study (Section~\ref{sec:setup}).}
% Provenance: experiments/archive/39_label_free_proxy_index_train_val_test_over_judge_families/scripts/e39.py (MODELS, PROMPTS) and results.md ("Selectors actually used").
The six judge models were served through one hosted gateway under the identifiers \texttt{openai/gpt-6-sol-high} (reasoning effort high), \texttt{anthropic/claude-sonnet-5-medium} (compute effort medium), \texttt{openai/gpt-6-luna-high} (high), \texttt{google/gemini-3.1-pro-preview}, \texttt{anthropic/claude-opus-5-5-no-fallback} (compute effort maximum), and \texttt{alibaba/qwen3.6-27b}.
The agent judges ran the same models inside the Terminus2 harness (at most 20 turns, reasoning effort high where exposed) and wrote their verdict to a file.
The five fuller-context prompts are \textsc{VeriBench}'s frozen rubrics, imported byte-for-byte: \texttt{D0\_deployed\_tc\_v1} (integer $0$--$10$, divided by ten), \texttt{R4\_exp75\_human\_precondition}, \texttt{H0\_human\_mirror}, \texttt{H1\_human\_mirror\_rater\_conventions}, and \texttt{H2\_human\_mirror\_conventions\_demos} (integer $0$--$5$, divided by five); each shows the judge the complete reference, the task's Python source, and the full candidate.
The four legacy prompts are the \texttt{P1\_BASELINE}, \texttt{P4\_CHECKLIST}, \texttt{P8\_STRICT\_B}, and \texttt{BAD\_STRINGY} prompts listed below, verbatim, with the reference field holding the gold's last theorem block as in our earlier version.
Scores were normalized to $[0,1]$ at parse time, and the expert target is the $0$--$5$ mean divided by five, so no further scaling is applied.

\paragraph{Prompts of the earlier 72-pair validation set (Appendix~\ref{app:validation-set-72}).}
The primary family model identifiers are \texttt{gpt-5.3-codex}, \texttt{claude-opus-4-7}, and \texttt{Qwen/Qwen3-Coder-30B-A3B-Instruct}.
Each family panel contains the three semantic prompts below and \texttt{BAD\_STRINGY}, plus \texttt{gpt-4o-mini} with the checklist and strict prompts.
The other two configurations are the deterministic edit-distance control \texttt{editdist\_e1} and seeded control \texttt{random\_r1}.
All model prompts request scores in $[0,1]$, not integers on a $0$--$10$ scale.
The human ratings use $0$--$5$; correlations compare their ordering with judge scores without equating the scales.

The reference field retains its historical label ``REFERENCE IMPLEMENTATION,'' but the code fills it with the gold's final theorem block.
The candidate field contains the full candidate file for human-agreement scoring; property probes operate on the extracted theorem text.
If no theorem block is found, the extractor uses the full file.
The Codex panel aggregates three recorded calls for model configurations and nine for controls by the median.
The Opus and Qwen panels aggregate three calls for every configuration by the median.
These counts come from the saved repeat arrays; a summary-level repeat setting alone does not describe reused cached rows.
The random control is deterministic for its input and seed, so repeated identical scores do not imply a reliable semantic judge.
Prompt labels P1, P4, and P8 are historical identifiers and do not denote the property numbers P1--P4.

\noindent\textbf{\texttt{P1\_BASELINE}}\par
\vspace{0.25em}\begin{lstlisting}[style=promptstyle]
You are an expert in Lean 4 formal verification.
Score how well the GENERATED THEOREM captures the correctness properties
of the REFERENCE IMPLEMENTATION. Score from 0.0 to 1.0.

- 1.0: Semantically equivalent to what the reference guarantees
- 0.7-0.9: Captures most properties, misses minor ones
- 0.4-0.6: Captures some properties, has meaningful gaps
- 0.1-0.3: Technically valid Lean but misses core properties
- 0.0: Trivially true (e.g., `True`) or incorrect

REFERENCE IMPLEMENTATION:
{reference_impl}

GENERATED THEOREM:
{generated_theorem}

Respond ONLY with valid JSON, no markdown:
{"score": <float 0.0-1.0>,
 "reasoning": "<one sentence max>"}
\end{lstlisting}

\noindent\textbf{\texttt{P4\_CHECKLIST}}\par
\vspace{0.25em}
\begin{lstlisting}[style=promptstyle]
Evaluate this in order:
1) What properties does REFERENCE IMPLEMENTATION guarantee?
2) Which are preserved by GENERATED THEOREM?
3) Are there missing core guarantees or trivialization?
4) Map to score in [0.0, 1.0]:
   1.0 equivalent; 0.7-0.9 minor gaps;
   0.4-0.6 meaningful gaps; 0.1-0.3 weak;
   0.0 trivial/incorrect.

REFERENCE IMPLEMENTATION:
{reference_impl}

GENERATED THEOREM:
{generated_theorem}

Return JSON only:
{"score": <float 0.0-1.0>,
 "reasoning": "<one sentence>"}
\end{lstlisting}

\noindent\textbf{\texttt{P8\_STRICT\_B}}\par
\vspace{0.25em}
\begin{lstlisting}[style=promptstyle]
Evaluate conservatively. If uncertain, choose the lower score.
Use [0.0,1.0]:
1.0 equivalent semantics
0.7-0.9 mostly right, minor misses
0.4-0.6 substantial omissions
0.1-0.3 weak relation
0.0 trivial/incorrect

REFERENCE IMPLEMENTATION:
{reference_impl}

GENERATED THEOREM:
{generated_theorem}

Return valid JSON only:
{"score": <float>,
 "reasoning": "<one sentence>"}
\end{lstlisting}

\noindent\textbf{\texttt{BAD\_STRINGY negative control}}\par
\vspace{0.25em}
\begin{lstlisting}[style=promptstyle]
You are a text similarity grader.
Compare generated theorem text to the reference implementation text.
Do NOT evaluate semantic correctness or theorem meaning.
Score ONLY based on superficial lexical/format overlap:
- high score if tokens/phrases look similar
- low score if wording/structure differs
Reference implementation:
{reference_impl}

Generated theorem:
{generated_theorem}

Return ONLY a JSON object:
{"score": <float 0.0 to 1.0>,
 "reasoning": "<short reason>"}
\end{lstlisting}


